%=====================================================================
\chapter{Shortest Paths: Single-Source and All-Pairs}\label{ch:shortestpaths}
%=====================================================================
\locator{5}{Part V --- Graph Algorithms}

\begin{outcomes}
By the end of this chapter you will be able to:
\begin{tight}
  \item \textbf{Solve} single-source shortest-path problems with Dijkstra's and
        Bellman--Ford's algorithms \emph{(CLO-7)}.
  \item \textbf{Prove} Dijkstra's algorithm correct and \textbf{identify} the
        step that requires non-negative weights.
  \item \textbf{Construct} a graph on which Dijkstra's algorithm returns a
        wrong answer, and \textbf{explain} why.
  \item \textbf{Apply} Floyd--Warshall to the all-pairs problem and
        \textbf{justify} the loop order.
  \item \textbf{Choose} among the three algorithms given density, weights and
        the number of sources required.
\end{tight}
\end{outcomes}

\begin{prereq}
\chref{graphs} for representation and BFS --- Dijkstra is BFS with a priority
queue. \chref{greedy} for the exchange argument. \chref{dpone} for
Floyd--Warshall, which is a dynamic program.
\end{prereq}

\begin{keyterms}
shortest path $\bullet$ relaxation $\bullet$ Dijkstra $\bullet$ Bellman--Ford
$\bullet$ Floyd--Warshall $\bullet$ negative edge $\bullet$ negative cycle
$\bullet$ single-source $\bullet$ all-pairs $\bullet$ optimal substructure
$\bullet$ early termination
\end{keyterms}

\section{Relaxation, the Common Operation}

\begin{definitionbox}[Relaxation]
Each algorithm here maintains an estimate $d[v] \ge \delta(s,v)$, the true
shortest distance. \term{Relaxing} an edge $(u,v)$ asks whether going via $u$
improves it:
\[ \textbf{if } d[u] + w(u,v) < d[v] \textbf{ then } d[v] \gets d[u]+w(u,v). \]

\smallskip
All three algorithms do only this. They differ entirely in \emph{the order} in
which edges are relaxed, and that order is what buys or costs the running time.
\end{definitionbox}

\begin{conceptbox}[Optimal substructure, and where it comes from]
Any subpath of a shortest path is itself a shortest path: if $p$ is a shortest
path from $s$ to $v$ passing through $u$, the portion from $s$ to $u$ must be
shortest, or substituting a shorter one would improve $p$.

\smallskip
That is a cut-and-paste argument, and it is the optimal substructure of
\chref{dpone}. It is why shortest paths yield to both greedy (Dijkstra) and
dynamic programming (Floyd--Warshall) --- and why \emph{longest simple paths}
do not, since a subpath of a longest simple path need not be a longest simple
path.
\end{conceptbox}

\section{Dijkstra}

\dsfig{\aoaalg{\algname{Dijkstra}$(G, w, s)$}{
\For{each vertex $v$}
  \State $d[v] \gets \infty$
\EndFor
\State $d[s] \gets 0$;\quad $S \gets \emptyset$;\quad
       $Q \gets$ all vertices
\While{$Q$ is not empty}
  \State $u \gets \algname{Extract-Min}(Q)$
    \Comment{smallest estimate}
  \State $S \gets S \cup \{u\}$
    \Comment{$d[u]$ is now final}
  \For{each $v$ adjacent to $u$}
    \State relax $(u,v)$
  \EndFor
\EndWhile
}}{\algname{Dijkstra}. Line 7 is the greedy commitment: once a vertex is
extracted its distance is declared final and never reconsidered.
Theorem~\ref{thm:dijkstra} is what justifies that, and
Section~20.4 is what happens when its hypothesis fails.}{dijkstra}

\begin{theorem}[Dijkstra is correct]\label{thm:dijkstra}
If every edge weight is non-negative, then when \algname{Dijkstra} adds $u$ to
$S$, $d[u] = \delta(s,u)$.
\end{theorem}

\begin{proof}
Suppose not, and let $u$ be the first vertex added to $S$ with
$d[u] \ne \delta(s,u)$. Then $u \ne s$, and $\delta(s,u) < \infty$ (otherwise
$d[u] = \infty = \delta(s,u)$).

\smallskip
Consider a shortest path $p$ from $s$ to $u$. Since $s \in S$ and $u \notin S$
at the moment before $u$ is added, $p$ must leave $S$ at some point: let $y$ be
the first vertex on $p$ not in $S$, and $x \in S$ its predecessor.

\smallskip
Because $x$ was added to $S$ earlier and $u$ is the \emph{first} mistake,
$d[x] = \delta(s,x)$. The edge $(x,y)$ was relaxed when $x$ was added, so
\[ d[y] \;=\; \delta(s,x) + w(x,y) \;=\; \delta(s,y), \]
the second equality because the prefix of a shortest path is a shortest path.

\smallskip
Now $y$ appears on a shortest path to $u$ and \textbf{all remaining edge
weights are non-negative}, so
\[ \delta(s,y) \;\le\; \delta(s,u). \]
But $u$ was chosen by \algname{Extract-Min} over $y$, so $d[u] \le d[y]$. Chain
them:
\[ d[u] \;\le\; d[y] \;=\; \delta(s,y) \;\le\; \delta(s,u) \;\le\; d[u], \]
the last because $d$ never underestimates. So all are equal and
$d[u] = \delta(s,u)$ --- contradicting the choice of $u$.
\end{proof}

\begin{alertbox}[Exactly one step used non-negativity]
The inequality $\delta(s,y) \le \delta(s,u)$. It says that extending a path can
never make it cheaper.

\smallskip
With a negative edge that is false: the rest of the path from $y$ to $u$ might
have negative total weight, so $\delta(s,u)$ could be \emph{less} than
$\delta(s,y)$, and a vertex finalised early may later turn out to have been
reachable more cheaply. The greedy commitment on line 7 is then unjustified.

\smallskip
Section~20.4 builds the graph where this happens.
\end{alertbox}

\begin{worked}[Which Priority Queue, Measured]
$|V| = 4{,}000$, density varying. \texttt{code/ch20\_sp.cpp} compares a binary
heap with an unsorted array (extract-min by linear scan):

\rowcolors{2}{white}{lightbg}
\begin{center}
\begin{tabular}{@{}rrrrl@{}}
\toprule
\rowcolor{primary!15}
avg degree & $|E|$ & heap (ms) & array (ms) & winner \\
\midrule
   2 &     7{,}998 &  0.36 & 46.77 & heap \\
  32 &   127{,}970 &  1.97 & 53.39 & heap \\
 512 & 2{,}047{,}499 &  6.40 & 62.63 & heap \\
2000 & 7{,}998{,}035 & 15.32 & 63.47 & heap \\
\bottomrule
\end{tabular}
\end{center}

\smallskip
\textbf{The array column is flat} --- 47 to 63 ms while $|E|$ grows 1000-fold
--- which is $\bigO{|V|^{2}}$ behaving exactly as advertised, just as the
adjacency matrix did in \chref{graphs}.

\smallskip
\textbf{But the textbook crossover does not appear.} The standard claim is that
the $\bigO{|V|^{2}}$ array version wins on dense graphs, since $|E| \approx
|V|^{2}$ makes the heap's $\lg|V|$ factor pure overhead. Here the heap won at
every density, including $|E| = 8 \times 10^{6}$ --- half of complete.

\smallskip
\textbf{The honest reading} is that at $|V| = 4{,}000$ the array's $|V|^{2} =
1.6\times10^{7}$ branch-laden scans cost more than $|E|\lg|V| = 10^{8}$
well-predicted heap operations on this processor. \chref{graphs} found the same
thing and fixed it by bit-packing; there is no equivalent trick here, because
the array version must examine real distances rather than single bits. The
crossover is real asymptotically and did not arrive at any size tested.
\end{worked}

\section{Bellman--Ford}

\dsfig{\aoaalg{\algname{Bellman-Ford}$(G, w, s)$}{
\For{each vertex $v$}
  \State $d[v] \gets \infty$
\EndFor
\State $d[s] \gets 0$
\For{$i \gets 1$ \textbf{to} $|V| - 1$}
  \For{each edge $(u,v)$}
    \State relax $(u,v)$
  \EndFor
\EndFor
\For{each edge $(u,v)$}
  \If{$d[u] + w(u,v) < d[v]$}
    \State \Return \textsc{negative-cycle}
  \EndIf
\EndFor
\State \Return $d$
}}{\algname{Bellman-Ford}. No priority queue and no assumption about signs:
$\bigO{|V||E|}$. The final pass is not decoration --- it is the only way to
detect a negative cycle, which Dijkstra cannot do at all.}{bellmanford}

\begin{theorem}[Bellman--Ford]\label{thm:bellman}
If $G$ has no negative-weight cycle reachable from $s$, then after $|V|-1$
passes $d[v] = \delta(s,v)$ for every $v$.
\end{theorem}

\begin{proof}
A shortest path has at most $|V|-1$ edges: if it had more it would repeat a
vertex, and the cycle so formed has non-negative weight (by hypothesis) and can
be removed without lengthening the path.

\smallskip
Now argue by induction that after $i$ passes, $d[v] = \delta(s,v)$ for every
$v$ whose shortest path uses at most $i$ edges. After 0 passes this holds for
$v = s$. For the step: let $v$ have a shortest path $s \leadsto u \to v$ using
$i+1$ edges; its prefix to $u$ uses $i$ edges and is shortest, so by hypothesis
$d[u] = \delta(s,u)$ after $i$ passes. Pass $i+1$ relaxes $(u,v)$ and sets
$d[v] \le \delta(s,u)+w(u,v) = \delta(s,v)$; since $d$ never underestimates,
equality holds.

\smallskip
After $|V|-1$ passes every shortest path is covered. If a further relaxation
still improves some $d[v]$, then some path uses $|V|$ or more edges and is
still improving, which requires a negative cycle.
\end{proof}

\begin{worked}[Bellman--Ford Is Faster Than Its Bound Suggests]
On random graphs with non-negative weights, average degree 8:

\rowcolors{2}{white}{lightbg}
\begin{center}
\begin{tabular}{@{}rrrrl@{}}
\toprule
\rowcolor{primary!15}
$|V|$ & Dijkstra (ms) & Bellman--Ford (ms) & ratio & agree? \\
\midrule
 2{,}000 &  0.40 &  0.3 & 0.8 & yes \\
 8{,}000 &  1.89 &  1.7 & 0.9 & yes \\
32{,}000 & 11.21 & 11.0 & 1.0 & yes \\
\bottomrule
\end{tabular}
\end{center}

\smallskip
\textbf{Bellman--Ford matched Dijkstra}, despite a bound of $\bigO{|V||E|}$
against $\bigO{|E|\lg|V|}$ --- a predicted factor of $|V|/\lg|V| \approx
2{,}000$.

\smallskip
\textbf{The reason is the early exit:} the implementation stops when a pass
changes nothing, and on a random graph the distances settle after a handful of
passes rather than $|V|-1$. The $\bigO{|V||E|}$ bound is a worst case that
random graphs do not come close to.

\smallskip
\textbf{So does the bound ever materialise?} It needs a graph that forces every
pass --- a long path whose edges are relaxed in the wrong order, so information
advances one vertex per pass:

\rowcolors{2}{white}{lightbg}
\begin{center}
\begin{tabular}{@{}rrrrr@{}}
\toprule
\rowcolor{primary!15}
$|V|$ & Dijkstra (ms) & Bellman--Ford (ms) & ratio & passes used \\
\midrule
1{,}000 & 0.03 &  1.4 &  44.9 &   999 \\
2{,}000 & 0.03 &  5.5 & 190.8 & 1{,}999 \\
4{,}000 & 0.06 & 21.8 & 386.9 & 3{,}999 \\
8{,}000 & 0.11 & 87.2 & \textbf{789.5} & 7{,}999 \\
\bottomrule
\end{tabular}
\end{center}

\smallskip
\textbf{There it is: 789$\times$ slower, and the ratio doubling with $|V|$.}
The graph is not adversarial in its weights --- every edge costs 1. Only the
\emph{order} the adjacency lists are visited in is unfavourable.

\smallskip
\textbf{The lesson is about benchmarking, not about Bellman--Ford.} Measuring
only on random inputs would have shown the two algorithms as equals. It is the
third time this book has made that point --- \chref{quicksort}'s sorted input,
\chref{structures}'s sorted insertions, \chref{hashing}'s structured keys ---
and it will not be the last.
\end{worked}

\section{One Negative Edge}

\begin{worked}[Dijkstra, Measurably Wrong]
Four vertices. $0 \to 1$ costs 3, $0 \to 2$ costs 4, $2 \to 1$ costs $-2$, and
$1 \to 3$ costs 1. The true distances are $d(1) = 2$ (via vertex 2) and
therefore $d(3) = 3$.

\rowcolors{2}{white}{lightbg}
\begin{center}
\begin{tabular}{@{}rrrrrl@{}}
\toprule
\rowcolor{primary!15}
vertex & Dijkstra (array) & Dijkstra (heap) & Bellman--Ford & correct & \\
\midrule
0 & 0 & 0 & 0 & 0 & \\
1 & 2 & 2 & 2 & 2 & \\
2 & 4 & 4 & 4 & 4 & \\
3 & \textbf{4} & 3 & 3 & 3 & \textbf{wrong} \\
\bottomrule
\end{tabular}
\end{center}

\smallskip
\textbf{Trace it.} The array version pops vertex 1 at distance 3 and
\emph{finalises} it, relaxing $1 \to 3$ to give $d(3) = 4$. It then pops vertex
2 and lowers $d(1)$ to 2 --- but vertex 1 is already done, so its outgoing edge
is never re-examined and $d(3)$ stays at 4.

\smallskip
\textbf{Notice that $d(1)$ itself ends up correct.} The error is only visible
at vertex 3, because the wrong value had already \emph{propagated}. The
three-vertex example people usually reach for does not fail, for exactly this
reason: the wrong value is overwritten before anything depends on it. A
counter-example has to let the mistake escape.

\smallskip
\textbf{And the heap version got it right} --- which is an accident worth
understanding. The lazy heap implementation never marks anything final; it
re-pushes whenever an estimate improves. On this graph that lets it recover.
It is behaving like Bellman--Ford, not like Dijkstra, and on adversarial graphs
with negative edges it can take exponential time. \textbf{``Dijkstra is wrong
on negative edges'' is a statement about the algorithm; whether your
\emph{program} visibly fails depends on which implementation you wrote.}

\smallskip
\textbf{Negative cycles.} On this graph Bellman--Ford reported none, correctly.
Adding a cycle $0 \to 1 \to 2 \to 0$ of total weight $-1$, it reported one ---
also correctly. Dijkstra cannot detect a negative cycle at all, and with one
present ``shortest path'' is undefined, since going round again is always
cheaper.
\end{worked}

\section{All Pairs}

\dsfig{\aoaalg{\algname{Floyd-Warshall}$(W)$}{
\State $D \gets W$
\For{$k \gets 1$ \textbf{to} $n$}
  \Comment{$k$ \emph{must} be outermost}
  \For{$i \gets 1$ \textbf{to} $n$}
    \For{$j \gets 1$ \textbf{to} $n$}
      \State $D[i][j] \gets \min\bigl(D[i][j],\;
             D[i][k] + D[k][j]\bigr)$
    \EndFor
  \EndFor
\EndFor
\State \Return $D$
}}{\algname{Floyd-Warshall}. Five lines, $\bigTh{n^{3}}$ regardless of density,
and the smallest constant of any algorithm in this chapter --- three nested
loops over contiguous memory.}{floyd}

\begin{conceptbox}[Why $k$ must be outermost]
This is a dynamic program, and $k$ indexes the subproblem.

\smallskip
Let $d^{(k)}[i][j]$ be the shortest distance from $i$ to $j$ using only
vertices $1 \ldots k$ as \emph{intermediates}. Then
\[ d^{(k)}[i][j] = \min\bigl(d^{(k-1)}[i][j],\;
   d^{(k-1)}[i][k] + d^{(k-1)}[k][j]\bigr), \]
--- either the best path avoids $k$, or it goes through $k$ exactly once.

\smallskip
The $k$ loop therefore advances the subproblem, and $i$ and $j$ merely fill in
the table for the current one. Putting $k$ inside would compute something that
is not a shortest path for any well-defined subproblem.

\smallskip
(The single array works because $d^{(k)}[i][k] = d^{(k-1)}[i][k]$ --- a path to
$k$ gains nothing by using $k$ as an intermediate --- so the row and column
being read are unchanged by the current pass.)
\end{conceptbox}

\begin{worked}[All Pairs, Measured]
$|V| = 600$. Floyd--Warshall against running Dijkstra from every vertex:

\rowcolors{2}{white}{lightbg}
\begin{center}
\begin{tabular}{@{}rrrrl@{}}
\toprule
\rowcolor{primary!15}
avg degree & $|E|$ & Floyd--Warshall (ms) & $|V| \times$ Dijkstra (ms) & winner \\
\midrule
  2 &   1{,}198 &  35.8 &  19.3 & Dijkstra \\
  8 &   4{,}793 & 125.7 &  59.3 & Dijkstra \\
 32 &  19{,}170 & 156.6 & 112.1 & Dijkstra \\
128 &  76{,}667 & 170.7 & 204.6 & Floyd \\
300 & 179{,}711 & 170.3 & 299.8 & Floyd \\
\bottomrule
\end{tabular}
\end{center}

\smallskip
\textbf{The crossover is near average degree 64}, or a density of about 0.1.

\smallskip
\textbf{Floyd--Warshall's column is nearly flat} --- 36 to 171 ms --- because
$\bigTh{|V|^{3}}$ does not depend on $|E|$. (It is not perfectly flat because
the early rows have many $\infty$ entries, which the implementation skips.)

\smallskip
\textbf{Repeated Dijkstra grows with $|E|$}, as $\bigO{|V|(|V|+|E|)\lg|V|}$
requires.

\smallskip
\textbf{And the practical point.} For a dense graph, five lines of
Floyd--Warshall beat $|V|$ invocations of a far more sophisticated algorithm.
Its constant is tiny: three nested loops over a contiguous array, perfectly
predictable, perfectly prefetchable. This is the opposite of
\chref{structures}'s pointer-chasing penalty, and it is the same cause.
\end{worked}

\begin{conceptbox}[Choosing]
\rowcolors{2}{white}{lightbg}
\begin{longtable}{@{}L{3.4cm}L{3.0cm}L{3.0cm}L{4.4cm}@{}}
\toprule
\rowcolor{primary!15}
\textbf{Algorithm} & \textbf{Time} & \textbf{Weights} & \textbf{Use when} \\
\midrule
\endhead
BFS & $\bigTh{V+E}$ & unweighted & all weights equal \\
Dijkstra (heap) & $\bigO{(V{+}E)\lg V}$ & non-negative & one source, the
  common case \\
Bellman--Ford & $\bigO{VE}$ & any & negative edges, or you must detect a
  negative cycle \\
Floyd--Warshall & $\bigTh{V^{3}}$ & any (no neg.\ cycle) & all pairs on a
  dense graph \\
\bottomrule
\end{longtable}

\smallskip
If all weights are equal, use BFS --- a priority queue over identical keys is
pure overhead, and \chref{graphs}'s measured 0.17 ms against Dijkstra's 0.36 ms
on the same sparse graph is the price of not noticing.
\end{conceptbox}

\begin{pitfall}
\begin{tight}
  \item \textbf{Using Dijkstra with negative edges.} Measured wrong, and the
        failure can hide: $d(1)$ was right and $d(3)$ was not.
  \item \textbf{Testing for that bug on too small a graph.} The obvious
        three-vertex example self-corrects. The error must be allowed to
        propagate.
  \item \textbf{Assuming your Dijkstra will visibly fail.} A lazy heap version
        may silently behave like Bellman--Ford --- correct here, exponential
        elsewhere.
  \item \textbf{Benchmarking Bellman--Ford only on random graphs.} Measured
        equal to Dijkstra there and 789$\times$ slower on a path.
  \item \textbf{Putting $k$ anywhere but outermost in Floyd--Warshall.} It
        indexes the subproblem; the result is not a shortest path.
  \item \textbf{Running Dijkstra $|V|$ times on a dense graph.} Measured
        1.8$\times$ slower than five lines of Floyd--Warshall.
  \item \textbf{Using Dijkstra on an unweighted graph.} BFS is simpler and
        faster.
  \item \textbf{Asking for shortest paths with a negative cycle present.} The
        question has no answer; detect it and say so.
\end{tight}
\end{pitfall}

\begin{lab}[Three Algorithms, and One That Is Wrong]
Reproduces all five tables. Expect thirty minutes.

\begin{lstlisting}[language=C++]
// ch20_sp.cpp -- three shortest-path algorithms, and when each wins.
// Build: cl /O2 /EHsc /std:c++17 ch20_sp.cpp         (see Appendix A)

// --- 1. Dijkstra with a binary heap --------------------------------
// The "lazy" form: push a new entry on every improvement and skip
// stale ones when popped. O((V+E) lg V), simpler than decrease-key.
//
// The invariant is the whole correctness argument: when a vertex is
// popped with the smallest tentative distance, that distance is final.
// It holds only because every edge weight is >= 0.

// --- 2. Dijkstra with a linear scan --------------------------------
// Same algorithm; the priority queue is an unsorted array, so
// extract-min is O(V) and the total is O(V^2 + E).

// --- 3. Bellman-Ford: relax every edge, V-1 times ------------------
// No priority queue and no assumption about signs. The bound is
// O(VE) -- the price of handling negative edges and of being able to
// DETECT a negative cycle, which Dijkstra cannot.
        if (!changed) break;        // early exit: often much faster

// --- 4. Floyd-Warshall: all pairs, three nested loops --------------
// The k loop must be OUTERMOST. That is not a style choice: k indexes
// the subproblem ("paths using only intermediate vertices < k"), and
// putting it inside would compute something that is not a shortest
// path.

// --- 5. Heap against array, across densities -----------------------
// --- 6. Dijkstra against Bellman-Ford ------------------------------
// --- 6b. The graph that forces every pass --------------------------
// A single path, with the adjacency lists visited in an order that
// propagates exactly one vertex per pass. Nothing adversarial about
// the weights; only the ORDER matters.

// --- 7. One negative edge, and Dijkstra is wrong -------------------
// 0 -> 1 costs 3, 0 -> 2 costs 4, 2 -> 1 costs -2, 1 -> 3 costs 1.
// True distances: d(1) = 2 via vertex 2, so d(3) = 3.
//
// Array Dijkstra pops vertex 1 at distance 3 and finalises it,
// relaxing 1 -> 3 to give d(3) = 4. It then pops vertex 2 and lowers
// d(1) to 2 -- but vertex 1 is done, so its edge to 3 is never
// re-examined and d(3) stays at 4. The error has PROPAGATED, which is
// why this graph fails where the obvious three-vertex example does
// not: there the wrong value was overwritten before anything
// depended on it.

// --- 8. All pairs: Floyd-Warshall against V x Dijkstra -------------
\end{lstlisting}

\textbf{What to notice.} Part 7 prints both Dijkstras. They disagree, and the
one that agrees with Bellman--Ford is the one that is not really Dijkstra. An
experiment that had shown only the heap version would have concluded, wrongly,
that negative edges are harmless.

\smallskip
\textbf{Then try to break it.} Add a negative cycle to the graph in part 7 and
run the lazy heap Dijkstra on it. Predict first. It will not terminate in any
reasonable time --- each trip round the cycle lowers an estimate, so it
re-pushes for ever. Bellman--Ford detects the cycle in one extra pass and says
so. That is the practical difference between an algorithm that assumes
something and one that checks.
\end{lab}

\begin{checkpoint}
\begin{tightnum}
  \item Identify the single step in the proof of Theorem~\ref{thm:dijkstra}
        that uses non-negativity, and state what goes wrong without it.
  \item Why does Bellman--Ford need exactly $|V|-1$ passes, and what does the
        $|V|$-th pass detect \emph{(CLO-7)}?
  \item Measured, Bellman--Ford matched Dijkstra on random graphs and was
        789$\times$ slower on a path. What does this say about the value of a
        worst-case bound?
\end{tightnum}
\end{checkpoint}

\begin{chaptersummary}
\begin{tight}
  \item All three algorithms do nothing but relax edges; they differ only in the
        order, and the order is what costs or saves.
  \item Dijkstra is greedy: the extracted minimum is declared final. The proof
        uses non-negativity at exactly one step --- that extending a path cannot
        make it cheaper.
  \item Measured: a binary heap beat an unsorted array at every density tested,
        including half of complete. The array's flat 47--63 ms is
        $\bigO{|V|^{2}}$ visible; the textbook crossover did not arrive.
  \item Bellman--Ford is $\bigO{|V||E|}$, handles negative edges, and is the
        only one of the three that can detect a negative cycle.
  \item Measured: with early termination it matched Dijkstra on random graphs,
        and was 789$\times$ slower on a path that forces all $|V|-1$ passes.
        Benchmarking only on random inputs would have missed this entirely.
  \item Measured: array Dijkstra returned $d(3) = 4$ where the answer is 3,
        on a four-vertex graph with one negative edge --- and the error was
        visible only because it had propagated.
  \item Floyd--Warshall is $\bigTh{|V|^{3}}$ with a tiny constant; $k$ must be
        the outermost loop because it indexes the subproblem. Measured, it beat
        $|V|$ runs of Dijkstra above a density of about 0.1.
\end{tight}
\end{chaptersummary}

\begin{reviewq}
  \item \textbf{(MCQ)} Dijkstra's algorithm requires: (a)~a connected graph;
        (b)~non-negative edge weights; (c)~an acyclic graph; (d)~integer
        weights.
  \item \textbf{(MCQ)} Floyd--Warshall runs in: (a)~$\bigTh{V^{2}}$;
        (b)~$\bigTh{VE}$; (c)~$\bigTh{V^{3}}$; (d)~$\bigTh{E\lg V}$.
  \item \textbf{(Short)} State the relaxation operation and explain why all
        three algorithms are built from it \emph{(CLO-7)}.
  \item \textbf{(Short)} Give a four-vertex graph on which Dijkstra fails, and
        explain why a three-vertex one will not do.
  \item \textbf{(Short)} Why must $k$ be the outermost loop in
        Floyd--Warshall?
  \item \textbf{(Applied)} Prove that if a graph has no negative cycle, some
        shortest path between any two vertices uses at most $|V|-1$ edges, and
        explain how Bellman--Ford's final pass uses this.
  \item \textbf{(Applied)} A routing system has 2{,}000 nodes, 8{,}000 links,
        needs all-pairs distances, and some link costs are negative subsidies.
        Choose an algorithm, give the arithmetic from this chapter's
        measurements, and state what you would check about the input first.
\end{reviewq}
